% s09 The measure and weak clustering
\section{The measure and weak clustering}

This section builds the measure of King's printed state in the curvature sector and draws from the authors' earlier paper on the torus limits \cite{SymPhase} the consequences of reflection positivity: the ergodic decomposition, clustering, the unique joint vacuum on a component, the growth of the components' moments, and weak \OSf\ on almost every component. The joint field with the Wick square is the subject of the tube sections. The measure is built from the characteristic functional, not from the lattice laws, which converge only locally (\cref{rem:s09:S2K}).

\emph{Conventions.} $\cV\defeq\Lambda^2\R^3$ with its Euclidean fibre norm. For a real $u\in C^\infty_c(\R^3;\Lambda^2)$ put $\cZ^K(u)\defeq\ZKinf(\mathrm iu)$, the one-member functional at the coefficient $\mathrm i$ (\cref{thm:s05:CK}(b), with the total-source convention of \cref{thm:s06:S1K}). Fix an even $r\ge14$. $\cH^{-r}(\R^3;\Lambda^2)$ is the Hermite--Sobolev space, the dual of the completion of $\cS(\R^3;\Lambda^2)$ under $\norm u_r^2=\sum_{\alpha,c}\lambda_\alpha^{r}\ip{h_{\alpha,c}}u^2$, where $(h_{\alpha,c})_{\alpha\in\N^3,\,c\le3}$ is the Hermite basis of $L^2(\R^3;\Lambda^2)$ and $\lambda_\alpha=2|\alpha|+4$ are the eigenvalues of $-\Delta+|x|^2+1$; $\fournorm\cdot$ is the weighted $C^4$ norm of \cref{lem:s02:E}, and $\fournorm u\le C\norm u_{10}$ on $\cS$ \cite[lem:D:hermite(a)]{SymPhase}. $X\defeq H^{-r}_{\mathrm{loc}}(\R^3;\Lambda^2)$ carries the actions $(\tau_a\Phi)(u)=\Phi(\tau_{-a}u)$, $(g\Phi)(u)=\Phi(g^{-1}\cdot u)$ for $g\in\Oh$ (or $O(3)$), $(\theta_j\Phi)(u)=\Phi(\sigma_j\cdot u)$ for the coordinate reflections $\sigma_j$, with the signed transport of $2$-forms of \cref{def:s02:sources}, and the local $\sigma$-algebras $\Sigma(U)$ generated by $\Phi\mapsto\Phi(u)$, $u\in C^\infty_c(U;\Lambda^2)$ real; $\theta_{j,c}\defeq\tau_{ce_j}\theta_j\tau_{-ce_j}$ is the reflection in the plane $x_j=c$ and $\rho_{j,c*}u$ the reflected source (\cref{def:s07:reflection}). These are the space, actions and $\sigma$-algebras of the abstract part of \cite{SymPhase}, restricted from $\Lambda^2\oplus\R$ to the curvature summand. The \emph{lattice curvature field} of King's cube is the random distribution $\Phi^K_{k,J}:u\mapsto F_{\Lk}(u)$ under $\avK\cdot$. $J_0(u)$ is the least $J$ with $\supp u$ at $\ell^\infty$-distance $>1$ from $\R^3\setminus\CJ$.

\begin{theorem}[the measure, the local lattice laws, weak clustering]\label{thm:s09:S2K}
The limits below exist by \cref{hyp:s03:EC}.
\begin{enumerate}[label=\textup{(\alph*)},leftmargin=*]
\item \textup{(The measure; uses \cref{thm:s05:CK,cor:s06:TK} only.)} For every even $r\ge14$ there is exactly one Borel probability measure $\mu^{K,H}$ on $\cH^{-r}(\R^3;\Lambda^2)$ whose mixed moments on $C^\infty_c$ are the Schwinger functions $S^K_n$ of King's printed state. Its characteristic functional on $C^\infty_c$, indeed on $\cS$, is $\cZ^K$ (extended by continuity); its moments on $\cS$ are the continuous extensions of the $S^K_n$ and satisfy the bounds of \cref{thm:s05:CK}(c). Its image $\mu^K_S$ under the restriction map $\iota:\cH^{-r}\to X$ is the unique Borel probability measure on $X$ with characteristic functional $\cZ^K$ on $C^\infty_c$, and also the unique one whose mixed moments on $C^\infty_c$ are the $S^K_n$. \emph{Local convergence of the lattice laws:} for every real $\chi\in C^\infty_c(\R^3)$ and every box $Q$ containing $\supp\chi$ in its interior, the laws $\nu^K_{k,J,\chi}$ of the localised lattice fields $\chi\Phi^K_{k,J}$ on $H^{-r}_{\mathrm{per}}(Q;\Lambda^2)$ converge weakly, along the full sequence $k\to\infty$ at every fixed $J\ge J_0(\chi)$, to a law $\nu^K_{J,\chi}$, and $\nu^K_{J,\chi}\Rightarrow\nu^K_\chi$ as $J\to\infty$, where $\nu^K_\chi$ is the law of $\chi\Phi$ under $\mu^K_S$. Nothing is asserted about the laws of the unlocalised lattice fields at fixed $J$ (\cref{rem:s09:S2K}).
\item \textup{(Invariance and weak clustering; uses also \cref{thm:s06:S1K,lem:s08:OS2K}.)} $\mu^K_S$ satisfies the hypotheses (H1)--(H3) of the abstract part of \cite{SymPhase}, is invariant under $\Cub$ about every point, and the abstract theorems of \cite{SymPhase} apply to it with their proofs: $\mu^K_S=\int\mu_\xi\,d\nu(\xi)$ with components that are translation-invariant and ergodic, $\theta_j$-invariant, reflection positive about every coordinate plane and satisfy (H2) \cite[thm:D:decomp]{SymPhase}; clustering in the mean \cite[thm:D:mean]{SymPhase} and under diagonal separation \cite[thm:D:diag]{SymPhase}; the Osterwalder--Schrader reconstruction in each coordinate direction with a unique joint vacuum on every component \cite[prop:D:recon, prop:D:fixed(iii)]{SymPhase}; the growth bound $|S^\xi_N(u_1,\dots,u_N)|\le\sigma^\xi_N\prod_i\norm{u_i}_{r+2}$, $\sigma^\xi_N\le\alpha_\beta(\xi)(N!)^\beta$ for every $\beta>1$, for the components' Schwartz moments \cite[cor:D:schwinger(iv)]{SymPhase}; and the cubic-invariant decomposition into mixtures of $k(\zeta)\in\{1,2,3,6\}$ ergodic measures \cite[add:Ohdecomp]{SymPhase}. The same holds on $\cH^{-r}$ for $\mu^{K,H}$ \cite[lem:D:hermite(d)]{SymPhase}. $\nu$-almost every component satisfies \OSf\ in the weak form, ergodicity under the translation group $\R^3$ \cite[add:OS4(iii)]{SymPhase}.
\item \textup{(The rotation corollary.)} $\cZ^K$ is $O(3)$-invariant, $\cZ^K(R\cdot u)=\cZ^K(u)$ for all $R\in O(3)$ and real $u$, by \cref{thm:s06:O3} (the rotation theorem \cref{hyp:s03:RC}); the clauses below use only this invariance. Then for $\nu$-a.e.\ $\xi$: $\mu_\xi$ is invariant under $\Euc$; $\E_{\mu_\xi}[\overline{F\circ\theta_{n,c}}\,F]\ge0$ for every unit vector $n$, every $c\in\R$ and every $F\in L^2(\Sigma(H^+_{n,c}),\mu_\xi)$, $H^+_{n,c}=\{x\cdot n>c\}$; and $\E_{\mu_\xi}[G'\,U(a)G]\to\E_{\mu_\xi}[G']\,\E_{\mu_\xi}[G]$ as $|a|\to\infty$ in any manner, for all $G,G'\in L^2(\mu_\xi)$. Consequently $\mu_\xi$ satisfies \OSf\ in each of the forms of \cite[def:OS4, add:OS4]{SymPhase} (the cluster property on the positive-time algebra, ergodicity of the time translations, ergodicity under $\R^3$) for every choice of the time direction, coordinate or not, and in the reconstruction from $\mu_\xi$ in any direction $n$ the joint fixed space is $\ker H_n=\C\Omega_n$ \cite[p:cor:rotcor]{SymPhase}. The moments of $\mu_\xi$ are real and symmetric; the vanishing of its odd curvature moments is \emph{not} claimed for a component.
\end{enumerate}
\end{theorem}

\subsection{Proof of (a)}

\emph{Step 1 (moment bounds, and where they hold).} For real $2$-forms $f_1,\dots,f_n$ at distance $>1$ from $\R^3\setminus C_{J_0}$, \cref{thm:s05:CK}(c) gives, at every $k$ and $J\ge J_0$, $|S^{K,(k,J)}_n|\le(\sqrt2\,e)^n(n!)^{1/2}\prod_j\norm{f_j}_{2,k}$ with $\norm f_{2,k}^2\le3C_0\fournorm f^2$ uniformly in $k$ (\cref{lem:s02:E}), and in the limit the bound of \cref{thm:s05:CK}(c); by \cref{thm:s05:CK}(a), $\avK{F(u)^2}\le2\norm u^2_{2,k}$ and hence $S^K_2(u,u)\le2\norm u_2^2$. The lattice Schwinger functions are the Taylor coefficients at $0$ of $a^K_{k,J}$, that is the moments of $\Phi^K_{k,J}$, and converge to the $S^K_n$. These bounds are uniform in $k$ and $J\ge J_0$ for sources at distance $>1$ inside $C_{J_0}$; at a fixed $J$ they are not available for sources reaching $\partial\CJ$, and neither is the convergence in $k$: \cref{hyp:s03:EC} and \cref{thm:s05:CK}(a) are stated for sources whose plaquettes have their corners in $\XJ$.

\emph{Step 2 (the measure).} (i) \emph{The functional on $\cS$.} For real $u_1,\dots,u_N\in C^\infty_c$ and $c\in\C^N$, at every $(k,J)$ with $J\ge J_0(u)$, $\sum_{l,m}\bar c_lc_m\,Z^K_{k,J}(\mathrm i(u_m-u_l))=\avK{|\sum_mc_me^{\mathrm iF(u_m)}|^2}\ge0$ by the linearity of $F_{\Lk}$ in the source. The limits $k\to\infty$, then $J\to\infty$, exist at these points by \cref{thm:s05:CK}(b) for the family $u$, so $\cZ^K$ is positive definite on $C^\infty_c$ with $\cZ^K(0)=1$. By \cref{cor:s06:TK}'s Lipschitz bound and (E2), $|\cZ^K(u)-\cZ^K(u')|\le\sqrt2\norm{u-u'}_2\le\frac\pi2\fournorm{u-u'}$. Since $\fournorm\cdot$ is a continuous seminorm on $\cS(\R^3;\Lambda^2)$ and $C^\infty_c$ is dense in $\cS$, $\cZ^K$ extends uniquely to a continuous function on $\cS$, still positive definite with $\cZ^K(0)=1$. By the Bochner--Minlos theorem \cite{GJ,Jan97} ($\cS$ is nuclear) there is exactly one Borel probability measure $\mu'$ on $\cS'(\R^3;\Lambda^2)$, with its cylinder $\sigma$-algebra, such that $\int e^{\mathrm i\Phi(u)}\,d\mu'=\cZ^K(u)$ for all real $u\in\cS$.

(ii) \emph{Moments on $C^\infty_c$, and on $\cS$ by uniform integrability.} For real $u_1,\dots,u_N\in C^\infty_c$ the characteristic function of $(\Phi(u_l))_l$ under $\mu'$ is $\varphi(t)=\cZ^K(\sum_lt_lu_l)=Z^{K,u}_\infty(\mathrm it)$, the restriction to $\mathrm i\R^N$ of the entire function $Z^{K,u}_\infty$ (\cref{thm:s05:CK}(b)). A characteristic function real-analytic near $0$ has moments of all orders, with $\int\prod_l\Phi(u_l)\,d\mu'=(-\mathrm i)^N\partial_{t_1}\cdots\partial_{t_N}\varphi(0)$ \cite{Luk70}; the chain rule gives $\partial_{t_1}\cdots\partial_{t_N}\varphi(0)=\mathrm i^N\partial_{z_1}\cdots\partial_{z_N}Z^{K,u}_\infty(0)$, so the factors cancel and
\begin{equation}\label{eq:s09:moments}
\int\prod_{l=1}^N\Phi(u_l)\,d\mu'=S^K_N(u_1,\dots,u_N);\qquad\int\Phi(u)^2\,d\mu'=S^K_2(u,u)\le2\norm u_2^2\le\tfrac{\pi^2}4\fournorm u^2\le C_2''\norm u_{10}^2
\end{equation}
by Step~1, (E2) and $\fournorm u\le C\norm u_{10}$. \emph{Extension to $\cS$.} By \cref{thm:s05:CK}(c) each $S^K_N$ extends to a continuous symmetric multilinear form on $\cS^N$ with $|S^K_N|\le\sigma_N\prod_l\fournorm{u_l}$. For $u_l\in\cS$ take cut-offs $\chi_R=\chi(\cdot/R)$, $\chi\in C^\infty_c$, $\chi=1$ on the unit ball, $0\le\chi\le1$, $R\ge1$: $\chi_Ru_l\to u_l$ in $\cS$, so $\Phi(\chi_Ru_l)\to\Phi(u_l)$ for every $\Phi\in\cS'$; and $\fournorm{\chi_Ru}\le C_\chi\fournorm u$ uniformly in $R\ge1$ with $C_\chi=2^4\norm\chi_{C^4}$, by Leibniz with $|R^{-|\beta|}(D^\beta\chi)(\cdot/R)|\le\norm\chi_{C^4}$ (the norm $\fournorm\cdot$ carries derivatives up to order $4$, so the cut-off does not pass through it with constant $1$). Hence $\int\prod_l\Phi(\chi_Ru_l)^2\,d\mu'=S^K_{2N}(\chi_Ru_1,\chi_Ru_1,\dots)\le\sigma_{2N}C_\chi^{2N}\prod_l\fournorm{u_l}^2$ uniformly in $R$: the products $\prod_l\Phi(\chi_Ru_l)$ are bounded in $L^2(\mu')$, hence uniformly integrable, and converge pointwise, so $\int\prod_l\Phi(u_l)\,d\mu'=\lim_RS^K_N(\chi_Ru_1,\dots)=S^K_N(u_1,\dots,u_N)$ by the continuity of $S^K_N$ on $\cS^N$.

(iii) \emph{The Hermite carrier.} Put $M(\Phi)\defeq\sum_{\alpha,c}\lambda_\alpha^{-r}\Phi(h_{\alpha,c})^2$. By \eqref{eq:s09:moments} with $\norm{h_{\alpha,c}}_{10}^2=\lambda_\alpha^{10}$, $\int M\,d\mu'\le3C_2''\sum_\alpha\lambda_\alpha^{10-r}<\infty$ as soon as $r-10>3$ (there are $(n+1)(n+2)/2$ multi-indices with $|\alpha|=n$), that is for every even $r\ge14$; so $\mu'(M<\infty)=1$. For $\Phi\in\cS'$ and $u\in\cS$ the Hermite expansion $u=\sum_{\alpha,c}\ip{h_{\alpha,c}}u\,h_{\alpha,c}$ converges in $\cS$ \cite[Thm.~V.13]{RS1}, so $\Phi(u)=\sum_{\alpha,c}\ip{h_{\alpha,c}}u\,\Phi(h_{\alpha,c})$ and, on $\{M<\infty\}$, $|\Phi(u)|\le M(\Phi)^{1/2}\norm u_r$ by Cauchy--Schwarz: $\Phi$ extends to $\Psi(\Phi)\in\cH^{-r}$ with $\norm{\Psi(\Phi)}_{-r}^2=M(\Phi)$ and $\Psi(\Phi)=\Phi$ on $\cS$. The map $\Phi\mapsto\Psi(\Phi)$ ($\defeq0$ off $\{M<\infty\}$) is measurable from the cylinder $\sigma$-algebra to the Borel $\sigma$-algebra of the separable Hilbert space $\cH^{-r}$, its coordinates being evaluations. Put $\mu^{K,H}\defeq\Psi_*\mu'$. Then $\int e^{\mathrm i\Psi(u)}\,d\mu^{K,H}=\cZ^K(u)$ and $\int\prod_l\Psi(u_l)\,d\mu^{K,H}=S^K_N(u_1,\dots,u_N)$ for all $u,u_l\in\cS$, by (ii). This is the carrier construction of \cite[lem:D:hermite(b)]{SymPhase} run on $\cS'$.

(iv) \emph{Uniqueness on $\cH^{-r}$.} By \cite[lem:D:hermite(c)]{SymPhase}, $\mu^{K,H}$ is the only Borel probability measure on $\cH^{-r}$ whose mixed moments on $C^\infty_c$ exist and equal the $S^K_n$. The proof there uses only that the moments $m_n$ of $Z=\sum_lt_l\Psi(u_l)$ satisfy $|m_n|\le c_S(n+1)K^nn!$ (then, by Cauchy--Schwarz for the odd absolute moments and the Taylor remainder bound $|y|^n/n!$, the characteristic function of $Z$ is the sum of its Taylor series on every interval of length $<2/K$, hence determined by the moments on all of $\R$) and the Lusin--Souslin theorem on $\cH^{-r}$ \cite{Kec95}. Our moments give the hypothesis with $c_S=1$: $m_n=S^K_n(g,\dots,g)$ for $g=\sum_lt_lu_l$, and $|m_n|\le(\sqrt2\,e)^n(n!)^{1/2}\norm g_2^n\le(n+1)K^nn!$ with $K=\sqrt2\,e\norm g_2$.

(v) \emph{The law on $X$ and its two uniqueness statements.} The map $\iota:\cH^{-r}\to X$, $\iota\Psi\defeq\Psi|_{C^\infty_c}$, is continuous, linear and injective \cite[lem:D:hermite(b)]{SymPhase}: $|\Psi(\chi f)|\le\norm\Psi_{-r}\norm{\chi f}_r\le C_\chi\norm\Psi_{-r}\norm f_{H^r}$ for a cut-off $\chi$, and $C^\infty_c$ is dense in $\cS$, which is dense in $\cH^r$. Put $\mu^K_S\defeq\iota_*\mu^{K,H}$; its characteristic functional on real $C^\infty_c$ is $\cZ^K$. By the uniqueness part of \cite[lem:D:realize]{SymPhase}, which applies to any Borel probability measure on $X$ (for a countable dense $(f_n)\subset C^\infty_c$ the sets $\{\Phi:(\Phi(f_1),\dots,\Phi(f_N))\in B\}$ form a $\pi$-system generating the Borel $\sigma$-algebra, by Lusin--Souslin, and the law of each such vector is determined by $t\mapsto\cZ^K(\sum_kt_kf_k)$), $\mu^K_S$ is the only Borel probability measure on $X$ with this characteristic functional. \emph{Moment uniqueness on $X$:} if $\mu''$ is any Borel probability measure on $X$ whose mixed moments on $C^\infty_c$ are the $S^K_n$, the argument of (iv) applies verbatim on $X$ with this $\pi$-system \cite[prop:discharge]{SymPhase}, so $\mu''=\mu^K_S$. And $\mu^{K,H}$ is the only measure on $\cH^{-r}$ with $\iota_*\mu^{K,H}=\mu^K_S$, since $\iota$ is a Borel isomorphism onto the Borel set $\iota\cH^{-r}$. Nothing in Step~2 sees the lattice: the measure clauses of (a) depend on King's printed state only through $\cZ^K$ and the $S^K_n$.

\emph{Step $2'$ (local convergence of the lattice laws).} Fix a real $\chi\in C^\infty_c(\R^3)$ and a box $Q$ with $\supp\chi$ compact in the interior of $Q$ (as in \cite[prop:C-tight]{SymPhase}), and let $(e_{\alpha,c})_{\alpha\in\Z^3,\,c\le3}$ be the Fourier basis of $L^2(Q;\Lambda^2)$, $|e_{\alpha,c}(x)|=|Q|^{-1/2}$. For a real $Q$-periodic smooth $2$-form $\psi$ put $\ip{\chi\Phi^K_{k,J}}\psi\defeq\Phi^K_{k,J}(\chi\psi)=F_{\Lk}(\chi\psi)$, where $\chi\psi\in C^\infty_c$ is supported in $\supp\chi$ (for the complex basis vectors, used only in (1), take real and imaginary parts); this identifies $\chi\Phi^K_{k,J}$ with an element of the real Hilbert space $H_\chi\defeq H^{-r}_{\mathrm{per}}(Q;\Lambda^2)$, and $\nu^K_{k,J,\chi}$ denotes its law.

(1) \emph{Second moments, uniformly.} For $J\ge J_0(\chi)$ and every $k$, every plaquette carrying $\chi e_{\alpha,c}$ has its centre in $\supp\chi$, hence its corners in $\XJ$, so the second-moment clause of \cref{thm:s05:CK}(a) applies to the real and imaginary parts: $\E^K_{k,J}|F_{\Lk}(\chi e_{\alpha,c})|^2\le2\norm{\chi\Re e_{\alpha,c}}^2_{2,k}+2\norm{\chi\Im e_{\alpha,c}}^2_{2,k}\le4\norm{\chi e_{\alpha,c}}^2_{2,k}$. \emph{Riemann-sum bound:} the plaquettes of one orientation have their centres on a translate of $\eps_k\Z^3$; those carrying $\chi e_{\alpha,c}$ have centres in $\supp\chi$; the cubes of side $\eps_k\le1$ about these centres are disjoint and lie in the $\frac12$-neighbourhood $N_{1/2}(\supp\chi)$; so $\norm{\chi e_{\alpha,c}}^2_{2,k}\le3\,\mathrm{vol}(N_{1/2}(\supp\chi))\,\norm\chi_\infty^2/|Q|=:K_\chi$, uniformly in $\alpha$, $c$, $k$ and $J\ge J_0(\chi)$ (no derivative of $\chi$ enters). Hence, for $s_0>\frac32$,
\begin{equation}\label{eq:s09:secondmoment}
\E^K_{k,J}\norm{\chi\Phi^K_{k,J}}^2_{H^{-s_0}_{\mathrm{per}}(Q)}=\sum_{\alpha,c}(1+|\alpha|^2)^{-s_0}\,\E^K_{k,J}|F_{\Lk}(\chi e_{\alpha,c})|^2\ \le\ 12K_\chi\sum_{\alpha\in\Z^3}(1+|\alpha|^2)^{-s_0}=:C'_\chi<\infty,
\end{equation}
uniformly in $k$ and $J\ge J_0(\chi)$ (the constants $4$ and $12K_\chi$ are valid but not sharp).

(2) \emph{Tightness.} The embedding $H^{-s_0}_{\mathrm{per}}(Q)\hookrightarrow H^{-r}_{\mathrm{per}}(Q)$ is compact for $r>s_0$ (Rellich on the torus); by Chebyshev, $\nu^K_{k,J,\chi}(\norm\cdot_{H^{-s_0}}>\rho)\le C'_\chi/\rho^2$, and the $H^{-s_0}$-ball of radius $\rho$ is compact in $H_\chi$; so $(\nu^K_{k,J,\chi})_{k\ge0,\,J\ge J_0(\chi)}$ is tight on the separable Hilbert space $H_\chi$, and Prokhorov's theorem \cite{Bil99} gives weak limit points along $k$ at fixed $J$.

(3) \emph{Identification at fixed $J$: the full sequence, and where \cref{hyp:s03:EC} enters.} For a real trigonometric polynomial $\psi$, $\E^K_{k,J}e^{\mathrm i\ip{\chi\Phi^K_{k,J}}\psi}=Z^K_{k,J}(\mathrm i\chi\psi)\to\ZKJ(\mathrm i\chi\psi)$ as $k\to\infty$, along the full sequence, by \cref{thm:s05:CK}(b) (that is \cref{cor:s04:G2}(ii)) for the one-member family $(\chi\psi)$ at $J\ge J_0(\chi)\ge J_0(\chi\psi)$: this is \cref{hyp:s03:EC} for that family on King's cube, together with the local bound of \cref{thm:s05:CK}(a) and Montel's theorem. A Borel probability measure on $H_\chi$ is determined by its characteristic functional on the dual $H^r_{\mathrm{per}}(Q;\Lambda^2)$, which is continuous there (dominated convergence), hence by its values on the dense set of real trigonometric polynomials. So every weak limit point of $(\nu^K_{k,J,\chi})_k$ is the one probability measure $\nu^K_{J,\chi}$ on $H_\chi$ with characteristic functional $\psi\mapsto\ZKJ(\mathrm i\chi\psi)$, and $\nu^K_{k,J,\chi}\Rightarrow\nu^K_{J,\chi}$ along the whole sequence $k\to\infty$ by tightness and the uniqueness of the limit point. Without \cref{hyp:s03:EC} the same argument gives the convergence along every continuum subsequence (\cref{thm:s12:LIM,cor:s12:KK}), not along the full sequence; the full-sequence clause is what \cref{hyp:s03:EC} adds for King's printed state.

(4) \emph{$J\to\infty$.} $\norm\cdot^2_{H^{-s_0}_{\mathrm{per}}}=\sum_{\alpha,c}(1+|\alpha|^2)^{-s_0}|\ip\cdot{e_{\alpha,c}}|^2$ is lower semicontinuous on $H_\chi$ (a supremum of finite sums of continuous functions), so by the portmanteau theorem \cite{Bil99} the bound \eqref{eq:s09:secondmoment} passes to $\nu^K_{J,\chi}$: $\int\norm\cdot^2_{H^{-s_0}}\,d\nu^K_{J,\chi}\le C'_\chi$ for every $J\ge J_0(\chi)$; hence $(\nu^K_{J,\chi})_J$ is tight on $H_\chi$ by the argument of (2), and $\ZKJ(\mathrm i\chi\psi)\to\ZKinf(\mathrm i\chi\psi)=\cZ^K(\chi\psi)$ (\cref{thm:s05:CK}(b)) identifies every limit point: $\nu^K_{J,\chi}\Rightarrow\nu^K_\chi$, the measure with characteristic functional $\psi\mapsto\cZ^K(\chi\psi)$.

(5) \emph{$\nu^K_\chi$ is the law of $\chi\Phi$ under $\mu^K_S$.} For $\Phi\in X$, $\chi\Phi\in H^{-r}(\R^3;\Lambda^2)$ is supported in $\supp\chi$, compact in the interior of $Q$, so its periodisation is an element of $H_\chi$ with $\ip{\chi\Phi}\psi=\Phi(\chi\psi)$ for $Q$-periodic smooth $\psi$; the map $\Phi\mapsto\chi\Phi$ is continuous from $X$ to $H_\chi$, since $|\Phi(\chi\psi)|\le\norm{\tilde\chi\Phi}_{H^{-r}}\norm{\chi\psi}_{H^r}\le c_\chi\norm{\tilde\chi\Phi}_{H^{-r}}\norm\psi_{H^r_{\mathrm{per}}(Q)}$ for a cut-off $\tilde\chi=1$ on $Q$. And $\int e^{\mathrm i\ip{\chi\Phi}\psi}\,d\mu^K_S=\int e^{\mathrm i\Phi(\chi\psi)}\,d\mu^K_S=\cZ^K(\chi\psi)$ by Step~2(v).

(6) \emph{Moments.} For every trigonometric polynomial $\psi$ (indeed every $\psi\in H^r_{\mathrm{per}}(Q;\Lambda^2)$, by the uniform-integrability argument of Step~2(ii) with $\chi\psi_n\to\chi\psi$), all moments of $\ip{\chi\Phi^K_{k,J}}\psi=F_{\Lk}(\chi\psi)$ are bounded uniformly in $k$ and $J\ge J_0(\chi)$ (Step~1; $\fournorm{\chi\psi}\le d_\chi\norm\psi_{C^4}$ by Leibniz), so by uniform integrability the moments of $\nu^K_{J,\chi}$ are the Taylor coefficients of $\ZKJ$ at $\chi\psi$, and those of $\nu^K_\chi$ are the $S^K_n(\chi\psi,\dots,\chi\psi)$, consistently with \eqref{eq:s09:moments}. This proves (a). \qed

\begin{remark}[why the convergence clause is local]\label{rem:s09:S2K}
The convergence $Z^K_{k,J}(\mathrm iu)\to\ZKJ(\mathrm iu)$ is available only for sources $u$ whose plaquettes have their corners in $\XJ$, that is for $u$ at distance $>1$ inside $\CJ$, where \cref{hyp:s03:EC} speaks. For a $u$ reaching $\partial\CJ$ the lattice functional is defined only after a convention for the plaquettes with a corner outside $\XJ$ (for instance, drop them), and for it no bound uniform in $k$ and no convergence is asserted: the second-moment bound of \cref{thm:s05:CK}(a) carries the corner hypothesis, and the crude bound $\avK{e^{F(\tilde u)}}\le e^{\frac12\lmu^{-2}\norm{\dstar_k\tilde u}^2_{2,k}}$ (\cref{lem:s02:GD} with $C\le\lmu^{-2}$ as forms) is not uniform in $k$ for a source $\tilde u$ truncated at $\partial\CJ$: $\dstar_k\tilde u$ is of order $\norm u_\infty/\eps_k$ on the $O(\eps_k^{-2}|\partial\CJ|)$ bonds of the truncation layer, so $\norm{\dstar_k\tilde u}^2_{2,k}=O(\eps_k^{-1})$. In particular the fixed-$J$ laws of the unlocalised lattice field on $X$ are not asserted to be tight; nothing below depends on them.
\end{remark}

\subsection{Proof of (b) and (c)}

\emph{Step 3 ((H1)--(H3)).} The hypotheses of the abstract part of \cite{SymPhase} are: (H1) a Polish space with a jointly continuous $\R^3$-action and three involutive homeomorphisms $\theta_j$ with $\theta_j\tau_a=\tau_{\rho_ja}\theta_j$, and $\mu$ invariant under all of them; (H2) local $\sigma$-algebras $\Sigma(U)$ with monotonicity, $\Sigma(U)=\bigvee_{K\Subset U}\Sigma(K)$, $\bigvee_K\Sigma(K)=\cB$, countable generation by bounded continuous functions, and covariance under $U(a)$ and $\Theta_j$; (H3) $\E_\mu[\overline{F\circ\theta_{j,c}}\,F]\ge0$ for $F\in L^2(\Sigma(H^+_{j,c}),\mu)$. On $X$ with the actions of the Conventions, the structural parts of (H1) and all of (H2) are proved in \cite[prop:discharge]{SymPhase} by arguments that use only the space and the push-forward actions (joint continuity of the translations on $H^{-s}_{\mathrm{loc}}$; the $\theta_j$ are linear homeomorphisms; Lusin--Souslin for $\bigvee_K\Sigma(K)=\cB$; countable dense families $\arctan\Phi(f_n)$; $\supp\tau_af=\supp f+a$, $\supp\sigma_j\cdot f=\sigma_j\supp f$); they hold verbatim for our $X$, that paper's space restricted to the curvature summand. \emph{Invariance of $\mu^K_S$:} by the uniqueness in Step~2(v) it suffices that $\cZ^K(\tau_au)=\cZ^K(u)$ and $\cZ^K(\sigma_j\cdot u)=\cZ^K(u)$ for real $u$; the first is \cref{cor:s06:TK} at the coefficient $\mathrm i$, the second \cref{thm:s06:S1K} at $z=\mathrm i$ with $O=\sigma_j$, and the same theorem with $O=\tau_bg\tau_{-b}$ gives the $\Oh$-invariance about every point. \emph{(H3):} for $E_a=e^{\mathrm i\Phi(u_a)}$ with $u_a$ real in $C^\infty_c(\{x_j>c\};\Lambda^2)$, $\overline{E_a\circ\theta_{j,c}}\,E_b=e^{\mathrm i\Phi(u_b-\rho_{j,c*}u_a)}$ and
\[
\E_{\mu^K_S}\big[\overline{E_a\circ\theta_{j,c}}\,E_b\big]=\cZ^K(u_b-\rho_{j,c*}u_a)=\ZKinf(\theta\zeta_a^*\sqcup\zeta_b),\qquad\zeta_a=(\mathrm i,u_a),
\]
the conjugated coefficient $-\mathrm i$ sitting on the reflected source as in \cref{def:s07:reflection}. By \cref{lem:s08:OS2K}, for the plane $x_j=c$ and every $c\in\R$, this Gram matrix is positive semidefinite. The exponential algebra is dense in $L^2(\Sigma(H^+_{j,c}),\mu^K_S)$ \cite[lem:A:mc]{SymPhase} (an algebra of bounded complex functions, closed under conjugation and containing the constants, is dense in $L^2$ of the $\sigma$-algebra it generates), and $\theta_{j,c}$ preserves $\mu^K_S$, so (H3) follows as in \cite[prop:discharge]{SymPhase}. Finally \cite[lem:D:hermite(d)]{SymPhase} carries (H1)--(H3) and the $\Oh$-invariance about every point from $(X,\mu^K_S)$ to $(\cH^{-r},\mu^{K,H})$.

\emph{Step 4 (weak clustering).} The standing hypothesis of the abstract part of \cite{SymPhase} is that $\mu$ is any measure satisfying (H1)--(H3); so \cite[thm:D:decomp]{SymPhase} gives the ergodic decomposition $\mu^K_S=\int\mu_\xi\,d\nu$ with components translation-invariant and ergodic, $\theta_j$-invariant, satisfying (H3) and (H2); \cite[thm:D:mean, thm:D:diag]{SymPhase} give clustering in the mean and under diagonal separation; \cite[prop:D:recon]{SymPhase} is the reconstruction in direction $j$, and on an ergodic component the joint fixed space is $\C\Omega_j$ \cite[prop:D:fixed(iii)]{SymPhase}. The proof of \cite[cor:D:schwinger(iv)]{SymPhase} uses only the $\Ezz$ bound in the form $\E_{\mu_S}Z_{\alpha,c}^{2n}\le c_SC_0^{2n}(2n)!\,\fournorm{h_{\alpha,c}}^{2n}$ on $Z_{\alpha,c}=\Phi(h_{\alpha,c})$, together with \cite[lem:D:hermite(a)]{SymPhase}, Jensen, Markov and Borel--Cantelli; \cref{thm:s05:CK}(c) supplies that bound on $\cS$ through \eqref{eq:s09:moments}, with $(e\pi/2)^{2n}((2n)!)^{1/2}\le(e\pi/2)^{2n}(2n)!$, so $c_S=1$, $C_0=e\pi/2$, and the growth bound of (b) holds for our components with the same proof. \cite[add:Ohdecomp]{SymPhase} assumes in addition that $\Oh$ acts on $X$ by homeomorphisms with $\rho_j=\theta_j$, $g\tau_ag^{-1}=\tau_{ga}$, $\mu$ invariant under every $g$, and $G\mapsto G\circ g^{-1}$ mapping $\Sigma(U)$-measurable functions to $\Sigma(gU)$-measurable ones: structural for the push-forward action ($(g\Phi)(u)=\Phi(g^{-1}\cdot u)$ is a linear homeomorphism of $X$, the relations are identities on test functions, $\supp g\cdot u=g\supp u$) together with the invariance of Step~3; it gives the cubic-invariant decomposition with $k(\zeta)\in\{1,2,3,6\}$. Weak \OSf\ is \cite[add:OS4(iii)]{SymPhase}, ergodicity for the action of $\R^3$, the ergodicity clause of the decomposition for $\nu$-a.e.\ component.

\emph{Step 5 (the rotation corollary).} \cite[p:cor:rotcor]{SymPhase} assumes (H1)--(H3) and that $O(3)$ acts on $X$ by homeomorphisms $x\mapsto Rx$, continuously in $R$ for each $x$, with $\rho_jx=\theta_jx$ for the coordinate reflections, $R\tau_aR^{-1}=\tau_{Ra}$, $G\mapsto G\circ R^{-1}$ mapping $\Sigma(U)$-measurable functions to $\Sigma(RU)$-measurable ones, and $\mu$ invariant under every $R\in O(3)$. The structural parts hold for $(R\Phi)(u)=\Phi(R^{-1}\cdot u)$ on $X$: the action is by linear homeomorphisms, continuous in $R$, and covariant for the local $\sigma$-algebras, as noted in \cite[add:rotcond-rem(1)]{SymPhase} for $H^{-s}_{\mathrm{loc}}(\R^3;\cV)$, and the push-forward respects the restriction to the $\Lambda^2$ summand. The invariance of $\mu^K_S$ is, by Step~2(v), $\cZ^K(R\cdot u)=\cZ^K(u)$ for real $u$, the hypothesis of (c), that is the $O(3)$-invariance of $\ZKinf$ at imaginary coefficients, supplied by \cref{thm:s06:O3}; this section does not derive it. The conclusions are those quoted in (c), proved there from the $R$-invariance and (H3) (reflection positivity about every plane), the argument of \cite[thm:D:decomp]{SymPhase} with a countable dense subgroup (the Euclidean invariance and reflection positivity of a.e.\ component), \cite[thm:D:diag]{SymPhase} along $(1,1,1)/\sqrt3$ after a rotation (clustering in every manner), and \cite[add:OS4rel(a)]{SymPhase} in a frame containing $n$ ($\ker H_n=\C\Omega_n$). They contain nothing on \OSth: the vanishing of the odd curvature moments of a component is not among them, and is not known for a single component.

\emph{Step 6 (curvature alone).} Everything above runs on $X=H^{-r}_{\mathrm{loc}}(\R^3;\Lambda^2)$ for the curvature field alone: the regularity input is \cref{thm:s05:CK}, the reflection positivity input \cref{lem:s08:OS2K}, the invariance input \cref{thm:s06:S1K}, the $O(3)$ input that of (c). The result is an ergodic decomposition of the law of the curvature field; nothing is claimed about the components of a joint law with the Wick square, and the single-axis statements of \cite{SymPhase} that rest on a torus-specific domination are not claimed for King's state. \qed

\begin{remark}[the hypotheses of the cited results, matched]\label{rem:s09:v1match}
For each result of \cite{SymPhase} used above, the step that supplies its printed hypotheses. (H1)--(H3): Step~3 (structure from prop:discharge; invariance from \cref{cor:s06:TK,thm:s06:S1K}; positivity from \cref{lem:s08:OS2K} and lem:A:mc). lem:D:realize, uniqueness part (a measure on $X$ with a given characteristic functional; Lusin--Souslin, lem:A:LS): Step~2(v). prop:discharge (the space and the push-forward actions; ``Determination by $S$''): Steps~3 and 2(v). lem:D:hermite: (a) no hypothesis; (b) the two-point bound $\E\Phi(f)^2\le C\norm f_{10}^2$, Step~2(ii); (c) $|m_n|\le c_S(n+1)K^nn!$, $r\ge14$ even, Step~2(iv) with $K=\sqrt2\,e\norm g_2$, $c_S=1$; (d) (H1)--(H3) on $X$, Step~3. The standing hypothesis ``$\mu$ is any measure satisfying (H1)--(H3)'' of thm:D:decomp, thm:D:mean, thm:D:diag, prop:D:recon, prop:D:fixed: Step~3; the ergodicity needed in prop:D:fixed(iii) holds on a component by thm:D:decomp. add:Ohdecomp (the $\Oh$ action and invariance): Step~4 and \cref{thm:s06:S1K}. cor:D:schwinger(iv) (the $\Ezz$ bound on Hermite functions; lem:D:hermite(a)): \cref{thm:s05:CK}(c) through Step~2(ii), $c_S=1$, $C_0=e\pi/2$. p:cor:rotcor (the $O(3)$ action and invariance): Step~5, the structure from add:rotcond-rem(1), the invariance from the hypothesis of (c). def:OS4, add:OS4(i)--(iii): definitions, Step~5; add:OS4(iii) is the ergodicity clause of thm:D:decomp. prop:C-tight, mechanism only (uniformly bounded second moments of $\Phi(\chi e_\alpha)$): Step~$2'$(1). lem:A:LS, lem:A:mc: Steps~2(iv),(v), 3. The results of \cite{SymPhase} stated for its own limit states (its clustering theorem, its Schwinger-function corollary as stated, its remark on components) are not cited; their abstract content is in the items above, and the odd-moment clause is not transferred to the components (clause (c)).
\end{remark}

\begin{remark}[the comparison family]\label{rem:s09:components}
The same conclusions hold for the limit states of the comparison family $\NDa$: for every continuum subsequence $\cK$ (\cref{def:s12:LIMstate}) the limit state $\ZLIM\cK$ carries a measure with the properties of \cref{thm:s09:S2K}(a) and (b), the local lattice laws converging along $\cK$ instead of the full sequence, and (c) (\cref{thm:s12:LIM}, clauses (v)--(vii)). By \cref{hyp:s03:EC} every such limit state equals $\ZKinf$ (\cref{cor:s04:G2}(iii)), so for King's printed state the two routes to the measure agree. For the iterated limit of the comparison family itself, when it is convergent at fixed cubes (\cref{def:s03:conv}), see \cref{prop:s12:ID}(b); nothing beyond \cref{cor:s04:G2}(iii) is asserted about it here, the two existence clauses not being interchangeable by the inputs of this paper (\cref{cor:s04:G2}(iv)).
\end{remark}
